\documentclass[a5j,12pt]{ltjtarticle}

\usepackage{hyperref}
\usepackage[noto-otf,match]{luatexja-preset}
\title{梅雨潭}
\author{朱自清}
\date{\today}

\begin{document}

\maketitle

\section{中国語・簡体字}

我第二次到仙岩的时候，我惊诧于梅雨潭的绿了。

梅雨潭是一个瀑布潭。仙瀑有三个瀑布，梅雨瀑最低。走到山边，便听见花花花花的声音；抬起头，镶在两条湿湿的黑边儿里的，一带白而发亮的水便呈现于眼前了。我们先到梅雨亭。梅雨亭正对着那条瀑布；坐在亭边，不必仰头，便可见它的全体了。亭下深深的便是梅雨潭。这个亭踞在突出的一角的岩石上，上下都空空儿的；仿佛一只苍鹰展着翼翅浮在天宇中一般。三面都是山，像半个环儿拥着；人如在井底了。这是一个秋季的薄阴的天气。微微的云在我们顶上流着；岩面与草丛都从润湿中透出几分油油的绿意。而瀑布也似乎分外的响了。那瀑布从上面冲下，仿佛已被扯成大小的几绺；不复是一幅整齐而平滑的布。岩上有许多棱角；瀑流经过时，作急剧的撞击，便飞花碎玉般乱溅着了。那溅着的水花，晶莹而多芒；远望去，像一朵朵小小的白梅，微雨似的纷纷落着。据说，这就是梅雨潭之所以得名了。但我觉得像杨花，格外确切些。轻风起来时，点点随风飘散，那更是杨花了。——这时偶然有几点送入我们温暖的怀里，便倏的钻了进去，再也寻它不着。





\end{document}
